% ID: MAT-GAN-ELL-005
% Titolo: Ellisse con vertice e retta tangente
% Disciplina: Matematica
% Area: geometria_analitica
% Argomento: ellisse
% Classe: Terza o quarta secondaria di secondo grado
% Difficolta: 4
% Tipo: problema_numerico
% Tag: matematica, geometria_analitica, ellisse
% Fonte: verifica_ellisse_anonimizzata
% Provenienza: materiale fornito dal docente; testo anonimizzato; soluzione redatta per Exergo

\begin{esercizio}
Trova l'equazione dell'ellisse che ha un vertice di coordinate $(2;0)$ ed è tangente alla retta di equazione
    $x-y-6=0.$
Assumi che il centro sia l'origine e gli assi siano paralleli agli assi cartesiani.
\end{esercizio}

\begin{soluzione}
Con le ipotesi esplicitate, $x^2/4+y^2/b^2=1$.
Sulla retta $y=x-6$, la sostituzione dà $(b^2+4)x^2-48x+144-4b^2=0$.
La tangenza impone discriminante $16b^2(32-b^2)=0$.
Escludendo $b=0$, risulta $b^2=32$ e l'equazione è $x^2/4+y^2/32=1$.
Senza centro e orientamento assegnati il problema originale non determina un'unica ellisse.
\end{soluzione}
